\chapter{Introdução}
\label{cap:introducao}

O processo de urbanização brasileiro consolidou, ao longo das últimas décadas, uma expressiva concentração demográfica nos centros urbanos. De acordo com os resultados do Censo Demográfico 2022, apurados pelo \citeonline{ibge2024censo2022malha}, cerca de 177,5 milhões de habitantes residem em áreas urbanas no Brasil, o que corresponde a 87,4\% da população total do país, estimada em 203,1 milhões de pessoas. Esse adensamento populacional exerce pressão contínua sobre os equipamentos e as redes físicas que sustentam a vida cotidiana nos municípios de médio e grande porte, exigindo da administração pública local ações permanentes de conservação, reparo e monitoramento do espaço urbano.

No exercício diário da gestão municipal, o conjunto de atividades voltadas à conservação das vias e dos equipamentos coletivos --- comumente denominado zeladoria urbana --- abrange serviços essenciais como a manutenção da pavimentação asfáltica, a operação da rede de iluminação pública, a desobstrução e o reparo de galerias de drenagem pluvial, o funcionamento da sinalização semafórica e a remoção de pontos de descarte irregular de resíduos sólidos. Quando essas infraestruturas sofrem degradação e não recebem intervenção tempestiva, instalam-se externalidades negativas imediatas que afetam a mobilidade, a segurança viária, a saúde pública e a valorização do patrimônio coletivo \cite{ibge2025entorno, snis2023drenagem}.

Apesar da relevância crítica desses serviços, o fluxo de comunicação e de gestão operacional predominante na maioria das prefeituras brasileiras ainda opera de maneira fragmentada e reativa. De um lado, o cidadão que identifica uma falha de infraestrutura em seu bairro encontra canais de atendimento dispersos, frequentemente desprovidos de padronização geográfica por Código de Endereçamento Postal (CEP) e incapazes de receber registros fotográficos ou audiovisuais que comprovem a gravidade da situação. De outro lado, os órgãos e secretarias municipais enfrentam dificuldades para consolidar demandas repetidas sobre um mesmo logradouro, planejar o deslocamento de suas equipes de manutenção e veículos de frota, controlar o consumo de materiais de almoxarifado adquiridos junto a fornecedores externos e aferir a efetiva resolução do problema junto ao solicitante.

A superação dessas lacunas informacionais relaciona-se diretamente aos compromissos globais pactuados na Agenda 2030 para o Desenvolvimento Sustentável da Organização das Nações Unidas \cite{onu2015agenda2030}. A estruturação de mecanismos transparentes, participativos e orientados por dados para a conservação do ambiente construído responde de forma central ao Objetivo de Desenvolvimento Sustentável (ODS) 11 (\emph{Cidades e Comunidades Sustentáveis}), especialmente no que tange à gestão participativa de assentamentos humanos e ao acesso seguro a espaços públicos qualificados \cite{onubrasil2015ods11, ipea2019ods11}. Adicionalmente, a modernização do controle operacional sobre ativos físicos e recursos públicos dialoga de maneira complementar com o ODS 9 (\emph{Indústria, Inovação e Infraestrutura}) e com o ODS 16 (\emph{Paz, Justiça e Instituições Eficazes}) \cite{onubrasil2015ods9, onubrasil2015ods16}.

Nesse contexto, a concepção de um sistema de banco de dados relacional apresenta-se como alicerce tecnológico indispensável para integrar, em uma fonte única e consistente, todas as etapas do ciclo de vida da manutenção urbana \cite{elmasri2019sistemas, heuser2009projeto}. O objetivo deste relatório, correspondente à Etapa 1 do projeto da disciplina de Banco de Dados, consiste em descrever em detalhes o problema do monitoramento e da manutenção colaborativa de infraestrutura urbana, fundamentar seu enquadramento nos ODS da ONU e delimitar o escopo funcional do sistema de banco de dados proposto, preparando a base conceitual para a modelagem entidade-relacionamento e lógica relacional da Etapa 2.

Para cumprir tais propósitos, este documento estrutura-se em cinco capítulos, incluindo esta introdução. O Capítulo~\ref{cap:problema} descreve em detalhes o problema de zeladoria e manutenção de infraestrutura urbana, seu contexto estatístico nacional, fluxo operacional, atores e lacunas informacionais. O Capítulo~\ref{cap:ods} analisa o enquadramento do tema nos ODS da Agenda 2030, examinando as metas e indicadores oficiais aplicáveis. O Capítulo~\ref{cap:escopo} define o escopo de alto nível do banco de dados, suas entidades de negócio, regras de integridade e consultas gerenciais em consonância com o modelo da Etapa 2. Por fim, o Capítulo~\ref{cap:consideracoes} apresenta as considerações finais e os encaminhamentos do projeto.
